\documentclass[11pt]{article}
\usepackage[margin=1in]{geometry}
\usepackage[T1]{fontenc}
\usepackage[utf8]{inputenc}
\usepackage{lmodern,microtype,amsmath,amssymb,amsthm,mathtools,booktabs}
\usepackage[colorlinks=true,linkcolor=blue,citecolor=blue,urlcolor=blue]{hyperref}
\hypersetup{pdftitle={Parity under standard-scaling neural-network Gaussian processes},pdfauthor={Claude}}
\setlength{\emergencystretch}{2em}
\newtheorem{theorem}{Theorem}[section]
\newtheorem{lemma}[theorem]{Lemma}
\newtheorem{proposition}[theorem]{Proposition}
\newtheorem{corollary}[theorem]{Corollary}
\theoremstyle{definition}\newtheorem{assumption}[theorem]{Assumption}
\theoremstyle{remark}\newtheorem{remark}[theorem]{Remark}
\DeclareMathOperator{\Var}{Var}
\DeclareMathOperator{\sign}{sign}
\DeclareMathOperator{\dist}{dist}
\DeclareMathOperator{\err}{err}
\DeclareMathOperator{\clip}{clip}
\newcommand{\E}{\mathbb E}
\newcommand{\Pp}{\mathbb P}
\newcommand{\X}{\mathcal X}
\newcommand{\R}{\mathbb R}
\newcommand{\sub}{\textup{[SE]}}
\title{Parity under standard-scaling neural-network\\Gaussian processes: superpolynomial sample\\lower bounds at fixed depth, and an obstruction}
\author{Claude}
\date{October 19, 2026}
\begin{document}
\maketitle
\begin{abstract}
We study Bayesian learning of parity by fully connected networks with the
standard neural-network Gaussian-process (NNGP) prior scaling: weight
variances $\sigma_1^2/(d+1)$ in the first layer and $\sigma_\ell^2/N$ in
later layers, at fixed depth, with width $N\to\infty$ first. The prior is
conditioned on an output margin $|z(x)|>a$ at every Boolean input and then
on the training signs; prediction counts as correct only when the correct
label has posterior probability greater than $2/3$. For $\tanh$ networks
of every fixed depth and every choice of layer gains, we prove that, for all large $d$,
$n\le\exp(\frac1{20}\sqrt d\log d)$ uniform independent examples of full
parity on $d$ bits leave, with probability $1-e^{-\Omega(d)}$ over the
sample, at least a $1/16$ fraction of inputs without confident correct
prediction. For the learner without the global margin, conditioned only on
the training signs, a threshold of the same form $d^{\Theta(\sqrt d)}$
(explicitly $d^{\sqrt d/10-1/2}/6000$ for $d\ge25$) holds on every
training set, for every bounded continuous nonconstant activation, every
depth and every gain; the argument is kernel-agnostic and also gives the
scale $d^k$ for fixed $k$-sparse parity. An $\exp(\Omega(d))$ lower bound,
the analogue of the subcritical theorem in this collection, remains open.
We prove that the step of that theorem which bounds the posterior's
conditioning cost by a polynomial per sample is false at the standard
scaling: for $\tanh$ networks, each fixed $\delta>0$ and
$\exp(d^{\epsilon})\le n\le e^{c(\delta)d}$, with probability at least
$1-2\delta$ the cost per sample is superpolynomial in $d$. We explain why
an upper bound on that cost matching our lower bound would complete an
exponential proof. Numerical spectra and small-dimensional
posterior computations illustrate the results.
\end{abstract}

\section{Introduction and results}
The collection's subcritical theorem \cite{subcritical} proves an
$\exp(\Omega(d))$ sample lower bound for full parity under a family of
Gaussian network priors whose width scaling is indexed by an exponent
$\gamma<1$. In that family the hidden preactivations diverge with width,
$\tanh$ acts as a sign function, and the limiting covariance is the
iterated arcsine kernel. The usual NNGP scaling is different: it keeps the
preactivations of order one, so the limit is the dual-activation kernel of
$\tanh$ itself. In the exponent notation of the collection's FAQ
(Question~9) this scaling is the tuple $(0,1,\ldots,1,1)$, which is not on
the one-parameter line $(\gamma-1,\gamma,\ldots,\gamma,\gamma+1)$.
The collection's NNGP note \cite{nngpnote} proves a $\Theta(2^d)$ bound
for a $\tanh$ NNGP whose depth grows linearly with $d$, and states that the
fixed-depth problem is open.

We ask whether parity needs exponentially many samples for the fixed-depth
standard-scaling learner. The answer we can prove is partial.

\paragraph{Results.} All statements take width to infinity first, at fixed
dimension, depth, gains, margin and training sample.
\begin{enumerate}
\item \textbf{Width limits} (Proposition~\ref{prop:width}, proved). For
every bounded continuous nonconstant activation the network output
converges to a nondegenerate Gaussian process with a dot-product kernel
$K=A\,F((1+x\cdot y)/(d+1))$, where $F$ has nonnegative Taylor coefficients
that do not depend on $d$. Both learners' posterior label probabilities
converge to those of the conditioned Gaussian process.
\item \textbf{Sign-only learner, every training set}
(Theorem~\ref{thm:signonly} and Corollary~\ref{cor:signonly}, proved).
Without the global margin, for every such activation, every depth and every
gain, every training set of size
$n\le d^{\sqrt d/10-1/2}/6000$ (with $d\ge25$; the explicit bound exceeds
one only for $d\ge385$) leaves, in the width limit, at least a quarter of the inputs of
full parity without confidence above $2/3$. For $k$-sparse parity the
threshold is $\binom d{\lfloor\sqrt k/5\rfloor}/6000$ (for $k\ge25$), and for fixed $k$
the failure fraction is at least $2^{-k}-O(n/\binom dk)$. The same bounds
hold for the margin learner when the readout variance is $\sigma_L^2N^{-b}$
with $b<1$ (Remark~\ref{rem:readout}).
\item \textbf{Margin learner, typical samples} (Theorem~\ref{thm:margin},
proved). For $\tanh$ at every fixed depth and every choice of gains and
margin, $n\le\exp(\frac1{20}\sqrt d\log d)$ iid uniform examples of full
parity leave, for $d\ge d_0$ (not explicit) and with probability at least $1-e^{-d/10}$, Gibbs error at least
$3/8$ and confidence failure at least $1/16$ in the width limit.
\item \textbf{The exponential question} (open). We do not know whether
$\exp(\Omega(d))$ samples are needed. Proposition~\ref{prop:obstruction}
(proved) shows that one specific step of the subcritical proof fails. That
step bounds the conditioning cost per sample by a polynomial in $d$ (at
most $Cd^{1+\beta}$). Here, for each fixed $\delta>0$, with probability at least $1-2\delta$ and for $\exp(d^\epsilon)\le n\le e^{c_1(\delta)d}$, that cost is superpolynomial in $d$. We
also identify a matching upper bound on that cost which would complete an
$\exp(\Omega(d))$ proof (Section~\ref{sec:obstruction}).
\end{enumerate}
The sign-only result uses only a universal inequality between consecutive
level masses of an isotropic kernel (Lemma~\ref{lem:levels}). The margin
result additionally needs a lower bound on every eigenvalue of the kernel.
We derive it from an explicit Fourier-integral formula for the Hermite
coefficients of $\tanh$ (Lemma~\ref{lem:tanh}).

\paragraph{What this paper does not show.} We prove no $\exp(\Omega(d))$
lower bound and no upper bound for either learner. We do not decide whether
the exponent $\sqrt d\log d$ is the truth. The margin theorem is proved for
$\tanh$, and for odd bounded continuous activations whose first-layer dual
coefficients are all positive and decay at most like
$e^{-o(\sqrt m\log m)}$, provided the linear coefficient $c_1(q_\ell)$
is positive at every later layer (Remark~\ref{rem:activations}). It holds for
$d\ge d_0$, where $d_0$ depends on the depth, gains and margin and is not
explicit, because we prove no rate for the coefficient tail $T(s)$. It does not cover
$\operatorname{erf}$, whose coefficients decay geometrically. As in
\cite{nngpnote,subcritical}, the results concern the exact posterior;
nothing is claimed about computing it.

\paragraph{Relation to kernel regression.} The NNGP posterior mean given
real-valued labels is kernel interpolation. For it, the span argument of
\cite[Proposition~4.1]{nngpnote} and the lower bounds of Hsu and of Simon
et al.\ for kernel regression concern squared error or correlation of the
mean. Our learners condition on signs, or on signs and a global margin, and
are judged by posterior confidence. That is a different criterion
(\cite[Section~5]{nngpnote}; FAQ Question~7).

\section{Model, width limits, and the cube spectrum}
\label{sec:model}
\subsection{Network and learners}
Let $\X=\{-1,1\}^d$ with the uniform probability measure, $M=2^d$ and
$D=d+1$. Fix a depth $L\ge2$, an activation $\phi:\R\to\R$, gains
$\sigma_1,\ldots,\sigma_L>0$, and a margin $a>0$. For width $N$ put
\[
 h^0(x)=(1,x)\in\R^D,\qquad h^\ell(x)=\phi(W^\ell h^{\ell-1}(x))
 \ (1\le\ell\le L-1),\qquad z_N(x)=W^Lh^{L-1}(x),
\]
with independent centered Gaussian weights of variances
\begin{equation}\label{eq:weights}
 \Var W^1_{ij}=\frac{\sigma_1^2}{D},\qquad
 \Var W^\ell_{ij}=\frac{\sigma_\ell^2}N\ (2\le\ell\le L-1),\qquad
 \Var W^L_i=\frac{\sigma_L^2}N .
\end{equation}
All hidden layers have width $N$, and there are no biases besides the
constant input coordinate. The \emph{unit-gain} network has every
$\sigma_\ell=1$. The specified-gain family of \cite[Section~7]{nngpnote}
has $\sigma_1=1$ and $\sigma_\ell^2=1/\E\tanh^2(G)$ for $\ell\ge2$, with
$G\sim N(0,1)$; at fixed depth it is another instance of
\eqref{eq:weights}.

\begin{assumption}\label{ass:act}
$\phi$ is bounded, continuous and not constant.
\end{assumption}

Let $f:\X\to\{-1,1\}$ be a target and $S=(X_1,\ldots,X_n)$ a sequence of
training inputs. Define the events
\begin{align*}
 B^{\rm mar}_S&=\{z:\ |z(x)|>a\ \forall x\in\X,\ f(X_i)z(X_i)>0\ \forall i\},\\
 B^{\rm sgn}_S&=\{z:\ f(X_i)z(X_i)>0\ \forall i\}.
\end{align*}
The \emph{margin learner} is the prior conditioned on $B^{\rm mar}_S$; this
is the collection's convention (FAQ Question~1). The \emph{sign-only
learner} is the prior conditioned on $B^{\rm sgn}_S$. For either learner put
\begin{align*}
 q_{N,S}(x)&=\Pp\{f(x)z_N(x)>0\mid z_N\in B_S\},\\
 F_N(S)&=\Pp_x\{q_{N,S}(x)\le2/3\},\qquad
 R_N(S)=\E_x[1-q_{N,S}(x)].
\end{align*}
$F_N$ is the confidence-failure fraction and $R_N$ the Gibbs error. For
every posterior and every $N$,
\begin{equation}\label{eq:mixture}
 R_N\le\frac{1-F_N}3+F_N,\qquad\text{so}\qquad F_N\ge\frac{3R_N-1}2,
\end{equation}
since $1-q\le1$ on the failure set and $1-q<1/3$ off it.

\subsection{Dual activations}
For $q>0$ write $v_\phi(q)=\E\phi(\sqrt qG)^2$ and, for $\rho\in[-1,1]$,
\[
 \widehat\phi_q(\rho)=\frac{\E\phi(\sqrt qU)\phi(\sqrt qV)}{v_\phi(q)},
\]
where $(U,V)$ are standard Gaussians with correlation $\rho$.
Under Assumption~\ref{ass:act}, $v_\phi(q)>0$.

\begin{lemma}[Dual activations]\label{lem:dual}
Expand $\phi(\sqrt q g)=\sum_{m\ge0}\beta_m He_m(g)/\sqrt{m!}$ in
$L^2(N(0,1))$, with probabilists' Hermite polynomials $He_m$. Then
\[
 \widehat\phi_q(\rho)=\sum_{m\ge0}c_m(q)\rho^m,\qquad
 c_m(q)=\frac{\beta_m^2}{\sum_k\beta_k^2}\ge0,\qquad \sum_mc_m(q)=1,
\]
with absolute convergence on $[-1,1]$. If $\phi$ is odd, $c_m(q)=0$ for
even $m$. Under Assumption~\ref{ass:act}, infinitely many $c_m(q)$ are
positive.
\end{lemma}
\begin{proof}
For correlated standard Gaussians,
$\E He_m(U)He_k(V)=\delta_{mk}\,m!\,\rho^m$; compare coefficients of
$s^mt^k$ in $\E e^{sU-s^2/2}e^{tV-t^2/2}=e^{\rho st}$. The bilinear
expansion of $\E\phi(\sqrt qU)\phi(\sqrt qV)$ then converges absolutely,
because $|\rho|\le1$ and $\sum\beta_m^2=v_\phi(q)<\infty$. Odd functions
have no even Hermite components. A bounded nonconstant function is not a
polynomial, so infinitely many $\beta_m$ are nonzero.
\end{proof}

Define $q_1=\sigma_1^2$, $v_\ell=v_\phi(q_\ell)$,
$q_{\ell+1}=\sigma_{\ell+1}^2v_\ell$ for $1\le\ell\le L-2$, and
\begin{equation}\label{eq:kernel}
 A=\sigma_L^2v_{L-1},\qquad
 F=\widehat\phi_{q_{L-1}}\circ\cdots\circ\widehat\phi_{q_1},\qquad
 K(x,y)=A\,F\!\left(\frac{1+x\cdot y}D\right).
\end{equation}
\begin{lemma}\label{lem:composite}
$F(t)=\sum_{m\ge0}c_mt^m$ with $c_m\ge0$ and $\sum_mc_m=1$. The
coefficients $c_m$ and the scale $A$ depend on $\phi$, $L$ and the gains,
but not on $d$. If $\phi$ is odd, $F$ is odd. Under
Assumption~\ref{ass:act}, infinitely many $c_m$ are positive.
\end{lemma}
\begin{proof}
Composition of power series with nonnegative coefficients summing to one
preserves both properties, and composition of odd series is odd. The
variances $q_\ell$ do not involve $d$, because $\|h^0(x)\|^2=D$ for every
input. For the last claim, induct on the number of layers. Suppose $H$ has
infinitely many positive coefficients and $H(1)=1$. Some $c_k(q)$ with
$k\ge1$ is positive; otherwise $\widehat\phi_q$ would be the constant
$c_0(q)=1$, forcing $\beta_m=0$ for $m\ge1$ and $\phi$ constant on the
support of a Gaussian. The series $H^{k-1}$ has some positive coefficient,
at degree $e$ say, so $H^k=H\cdot H^{k-1}$ has a positive coefficient at
degree $e+m$ for each of the infinitely many $m$ with $[t^m]H>0$.
\end{proof}

\subsection{Width limits}
\begin{proposition}[Width limits]\label{prop:width}
Under Assumption~\ref{ass:act}, fix $d,L$ and the gains.
\begin{enumerate}
\item[(i)] As $N\to\infty$, $z_N$ converges in law on $\R^\X$ to
$G\sim N(0,K)$, with $K$ as in \eqref{eq:kernel}.
\item[(ii)] $K$ is positive definite (Lemma~\ref{lem:spectrum}).
\item[(iii)] For every target $f$, training sequence $S$ and margin $a>0$,
and for either learner, $\Pp\{z_N\in B_S\}\to\Pp\{G\in B_S\}>0$, and
$q_{N,S}(x)\to q_S(x):=\Pp\{f(x)G(x)>0\mid G\in B_S\}$ for every $x$.
Consequently $R_N(S)\to R_S:=\E_x[1-q_S(x)]$,
$\liminf_NF_N(S)\ge(3R_S-1)/2$, and
$\liminf_NF_N(S)\ge\Pp_x\{q_S(x)<2/3\}$.
\end{enumerate}
\end{proposition}
\begin{proof}
(i) For a positive semidefinite matrix $\Sigma$ on $\X$ let
$\Gamma(\Sigma)=\E[\phi(u)\phi(u)^{\mathsf T}]$ with $u\sim N(0,\Sigma)$.
The map $\Gamma$ is continuous on the positive semidefinite cone: if
$\Sigma_k\to\Sigma$, then $\Sigma_k^{1/2}\xi\to\Sigma^{1/2}\xi$ for every
$\xi$, and dominated convergence applies because $\phi$ is bounded and
continuous. Write $Q^\ell_N(x,y)=N^{-1}\sum_ih^\ell_i(x)h^\ell_i(y)$ and
let $\mathcal F_\ell$ be generated by the first $\ell$ layers of weights.
The first-layer rows are independent, and each row's preactivation vector
is $N(0,\Sigma_1)$ with $\Sigma_1(x,y)=\sigma_1^2(1+x\cdot y)/D$.
Conditionally on $\mathcal F_{\ell-1}$ with $\ell\ge2$, the rows of the
$\ell$th layer are independent with preactivation law
$N(0,\sigma_\ell^2Q^{\ell-1}_N)$. Hence
$\E[Q^\ell_N\mid\mathcal F_{\ell-1}]=\Gamma(\sigma_\ell^2Q^{\ell-1}_N)$
(with $\Gamma(\Sigma_1)$ at $\ell=1$), and every entry has conditional
variance at most $\|\phi\|_\infty^4/N$. By Chebyshev's inequality and the
continuous mapping theorem, induction on $\ell$ gives
$Q^\ell_N\to K^\ell$ in probability, where $K^1=\Gamma(\Sigma_1)$ and
$K^\ell=\Gamma(\sigma_\ell^2K^{\ell-1})$. Every diagonal entry of
$\sigma_\ell^2K^{\ell-1}$ equals $q_\ell$, and the correlations follow
the recursion $\rho_1=(1+x\cdot y)/D$, $\rho_{\ell+1}=\widehat\phi_{q_\ell}
(\rho_\ell)$. Therefore $\sigma_L^2K^{L-1}=K$. Finally
$z_N\mid\mathcal F_{L-1}\sim N(0,\sigma_L^2Q^{L-1}_N)$, so for each
$t\in\R^\X$ the characteristic function is
$\E\exp(-\tfrac12\sigma_L^2t^{\mathsf T}Q^{L-1}_Nt)$, which converges to
$\exp(-\tfrac12t^{\mathsf T}Kt)$ by bounded convergence.

(iii) Every coordinate $G(x)$ has variance $A>0$. The boundary of
$B_S$, and of $B_S\cap\{f(x)z(x)>0\}$, lies in the finite union of
hyperplanes $\{z(y)=\pm a\}$ and $\{z(y)=0\}$, which is $G$-null. By the
portmanteau theorem these probabilities converge. The limit of the
denominator is positive because $K$ is nonsingular and $B_S$ is a nonempty
open set. The last assertions follow from \eqref{eq:mixture}, from the
finiteness of $\X$, and from the fact that $q_S(x)<2/3$ forces
$q_{N,S}(x)<2/3$ for all large $N$.
\end{proof}

The quantifier order is the one in the main results: $N\to\infty$ at fixed
$d$, $L$, gains, $a$ and $S$. No quantitative width is claimed.

\subsection{The kernel on the cube}
For $V\subseteq\{0,1,\ldots,d\}$ write $\chi_V(u)=\prod_{i\in V}u_i$ on
$\{-1,1\}^D$, and $\chi_U$ for $U\subseteq[d]$ on $\X$. Let $p_m(j)$ be
the probability that the walk on $\{-1,1\}^D$ which flips one uniformly
chosen coordinate per step is, after $m$ steps, at a specified vertex at
distance $j$ from its start.

\begin{lemma}[Cube spectrum; the kernel-free part of {\cite[Lemma~3.1]{subcritical}}]
\label{lem:spectrum}
Let $F=\sum_mc_mt^m$ with $c_m\ge0$ and $\sum c_m=1$. Then
$F(u\cdot u'/D)=\sum_Va_{|V|}\chi_V(u)\chi_V(u')$ on $\{-1,1\}^D$, with
$a_j=\sum_mc_mp_m(j)$, and:
\begin{enumerate}
\item[(a)] $a_j\ge a_{j+2}$ for $0\le j\le D-2$;
\item[(b)] $a_j>0$ if and only if $c_m>0$ for some $m\ge j$ with
$m\equiv j\pmod 2$;
\item[(c)] on $\X$, $K(x,y)=\sum_{U\subseteq[d]}\lambda_{|U|}\chi_U(x)
\chi_U(y)$ with $\lambda_q=A(a_q+a_{q+1})$, and
$\lambda_q$ is nonincreasing in $q$;
\item[(d)] if infinitely many $c_m$ are positive, every $\lambda_q$ is
positive and $K$ is positive definite;
\item[(e)] if $F$ is odd, $\lambda_{2k}=\lambda_{2k+1}=Aa_{2k+1}$.
\end{enumerate}
\end{lemma}
\begin{proof}
Expanding $(u\cdot u'/D)^m$ as an average over $m$ independent uniform
coordinates shows that its coefficient of $\chi_V(u)\chi_V(u')$ is
$p_m(|V|)$. For (a), the one-step recursion
$p_{m+1}(j)=\frac jDp_m(j-1)+\frac{D-j}Dp_m(j+1)$ gives
\[
 D[p_{m+1}(j)-p_{m+1}(j+2)]
 =j[p_m(j-1)-p_m(j+1)]+(D-j-2)[p_m(j+1)-p_m(j+3)],
\]
with $p_m(i):=0$ for $i\notin[0,D]$, and induction from $m=0$ gives
$p_m(j)\ge p_m(j+2)$ (the case $j+1=D-1$ of the hypothesis is
$p_m(D-1)\ge p_m(D+1)=0$, which is trivial); average against $c_m\ge0$. For (b),
$p_m(j)>0$ exactly when $m\ge j$ and $m\equiv j$. For (c), restrict to
$u=(1,x)$: $\chi_V(1,x)=\chi_{V\setminus\{0\}}(x)$, and each $U\subseteq[d]$
arises from $V=U$ and $V=U\cup\{0\}$. Monotonicity follows from (a):
$\lambda_{q+1}/A=a_{q+1}+a_{q+2}\le a_{q+1}+a_q$. For (d), given $q$ there
is $m\ge q+1$ with $c_m>0$, and by (b) one of $a_q,a_{q+1}$ (the one whose
index has the parity of $m$) is positive. The matrix $K$ has eigenvalue
$M\lambda_{|U|}$ on $\chi_U$. Part (e) is (c) with $a_{\rm even}=0$.
\end{proof}

Throughout, $\lambda_q$ denotes these \emph{normalized} eigenvalues of $K$
itself, so they include the factor $A$. For $g:\X\to\R$ write
$\widehat g(U)=\E_xg(x)\chi_U(x)$, $\|g\|_2^2=\E_xg(x)^2$,
$g_{\le r}=\sum_{|U|\le r}\widehat g(U)\chi_U$ and $g_{>r}=g-g_{\le r}$.
The Gaussian energy is $\|g\|_K^2=g^{\mathsf T}K^{-1}g=\sum_U\widehat g(U)^2/
\lambda_{|U|}$, the convention of \cite{subcritical}. Let
\[
 \tau_r=\sum_{q>r}\binom dq\lambda_q=\Var G_{>r}(x),\qquad
 T(r)=\sum_{m>r}c_m .
\]

\begin{lemma}[Level-mass inequalities]\label{lem:levels}
Let $F$ be as in Lemma~\ref{lem:spectrum} and $0\le r<d$.
\begin{enumerate}
\item[(a)] $\binom d{r+1}\lambda_{r+1}\le\tau_r$.
\item[(b)] If $F$ is odd and $r$ is odd, $\binom D{r+2}\lambda_{r+1}\le\tau_r$.
\item[(c)] $\tau_r\le A\,T(r)$.
\item[(d)] $s_{2,r}:=\sum_{q>r}\binom dq\lambda_q^2\le\lambda_{r+1}\tau_r$.
\end{enumerate}
\end{lemma}
\begin{proof}
(a) $\tau_r$ contains the term $q=r+1$. (b) Here $\lambda_{r+1}=\lambda_{r+2}$
by Lemma~\ref{lem:spectrum}(e), and
$\binom d{r+1}+\binom d{r+2}=\binom D{r+2}$. (c) By the grouping in
Lemma~\ref{lem:spectrum}(c),
$\tau_r/A=\sum_{V:\,|V\setminus\{0\}|>r}a_{|V|}\le\sum_{j>r}\binom Dja_j$,
and $\binom Dja_j=\sum_mc_m\Pp\{\text{walk at distance $j$ at time $m$}\}$.
Distance $j$ at time $m$ requires $m\ge j$, so
$\sum_{j>r}\binom Dja_j\le\sum_{m>r}c_m$. (d) Bound one factor
$\lambda_q$ by $\lambda_{r+1}$, using monotonicity.
\end{proof}

Inequality (a) is what makes the sign-only argument kernel-agnostic. The
subcritical paper uses the sharper arcsine ratio
$\lambda_{r+1}/\tau_r\asymp1/(r\binom D{r+2})$ \cite[Lemma~3.2]{subcritical};
we use only (a) and (b), which hold for every such kernel.

\section{The sign-only learner on every training set}
\label{sec:signonly}
In this section $K$ is any positive-definite covariance on $\X$ that
depends only on Hamming distance and whose degree eigenvalues
$\lambda_q$ are nonincreasing. By Lemma~\ref{lem:spectrum} this includes
every kernel \eqref{eq:kernel} under Assumption~\ref{ass:act}, for every
depth and gain. Let $G\sim N(0,K)$ and
\[
 q_S(x)=\Pp\{f(x)G(x)>0\mid f(X_i)G(X_i)>0\ \forall i\}.
\]
For $0\le r<d$ let
\[
 \err_r(f)=\min_{\deg p\le r}\Pp_x\{f(x)p(x)\le0\},
\]
the least error of a degree-$r$ polynomial threshold function with zeros
counted as errors. A small constant shift of $p$ removes its zeros without
creating new errors, so this agrees with the definition used in
\cite{beyondparity}.

\begin{lemma}[Gaussian mean and confidence {\cite[Lemma~10.1]{subcritical}}]
\label{lem:meanconf}
Let $\theta=13/20$ and $c_\theta=\sqrt{2\pi}(e^{-1}-7/20)$. Let
$m=\E[G\mid f(X_i)G(X_i)>0\ \forall i]$, let $s\le n$ be the number of
distinct training inputs, and let $\sigma_x^2=\Var(G(x)\mid G_S)$. Then for
every training sequence and every $0\le v<d$:
\begin{enumerate}
\item[(i)] $\|m\|_K^2\le s$;
\item[(ii)] if $\sigma_x>0$ and $q_S(x)>\theta$, then $f(x)m(x)>c_\theta\sigma_x$;
\item[(iii)] $\Pp_x\{\sigma_x^2<\tau_v/2\}\le2s\lambda_{v+1}/\tau_v\le2n\lambda_{v+1}/\tau_v$;
\item[(iv)] $\|m-m_{\le v}\|_2^2\le n\lambda_{v+1}$.
\end{enumerate}
\end{lemma}
\begin{proof}
This is \cite[Lemma~10.1]{subcritical}. Its proof uses Gaussianity,
positive definiteness, Pr\'ekopa's theorem and Fourier orthogonality.
It uses the kernel only through two facts: the high-degree part $G_{>v}$
is independent of $G_{\le v}$ and has constant variance $\tau_v$, and the
largest normalized eigenvalue above degree $v$ is $\lambda_{v+1}$. The
second fact holds here by monotonicity. For (iii): conditioning also on
$G_{\le v}$ can only lower the conditional variance, so
$\sigma_x^2\ge\tau_v-R(x,x)$, where
$R=Q_{\X S}Q_{SS}^{\dagger}Q_{S\X}$ and $Q$ is the covariance of $G_{>v}$.
$R$ is positive semidefinite of rank at most $s$, with
$\operatorname{tr}R\le s\|Q\|_{\rm op}=sM\lambda_{v+1}$. Markov's inequality
then gives (iii). Part (iv) follows from (i), since
$\|m_{>v}\|_2^2=\sum_{|U|>v}\widehat m(U)^2\le\lambda_{v+1}\|m\|_K^2$.
\end{proof}

\begin{theorem}[Sign-only learner]\label{thm:signonly}
For every target $f$, every sequence $S$ of $n$ inputs (repeats and
deliberately chosen inputs allowed), and every $0\le r<d$,
\begin{equation}\label{eq:signonly}
 \Pp_x\{q_S(x)\le13/20\}\ \ge\ \err_r(f)-C_0\frac{n\lambda_{r+1}}{\tau_r}
 \ \ge\ \err_r(f)-\frac{C_0\,n}{\binom d{r+1}},
 \qquad C_0=2+\frac2{c_\theta^2}<1000 .
\end{equation}
If the eigenvalues are paired as in Lemma~\ref{lem:spectrum}(e) and $r$ is
odd, $\binom d{r+1}$ may be replaced by $\binom{d+1}{r+2}$.
\end{theorem}
\begin{proof}
Let $p=m_{\le r}$. Let $E_1=\{\sigma_x^2<\tau_r/2\}$ and
$E_2=\{|m-p|\ge c_\theta\sqrt{\tau_r/2}\}$. By Lemma~\ref{lem:meanconf}(iii),
(iv) and Markov's inequality,
$\Pp(E_1)+\Pp(E_2)\le(2+2/c_\theta^2)n\lambda_{r+1}/\tau_r$. Outside
$E_1\cup E_2$, the inequality $q_S(x)>13/20$ implies
$f(x)m(x)>c_\theta\sigma_x\ge c_\theta\sqrt{\tau_r/2}>|m(x)-p(x)|$, hence
$f(x)p(x)>0$. Therefore
$\{fp\le0\}\subseteq\{q_S\le13/20\}\cup E_1\cup E_2$. Since $p$ has degree at
most $r$, $\Pp_x\{fp\le0\}\ge\err_r(f)$. The second inequality in
\eqref{eq:signonly} is Lemma~\ref{lem:levels}(a), and the paired version is
Lemma~\ref{lem:levels}(b). Numerically $c_\theta=0.04481\ldots$, so
$C_0=997.7\ldots$.
\end{proof}

\begin{lemma}[Parity is far from low-degree thresholds]\label{lem:abfr}
Let $\chi_k(x)=x_1\cdots x_k$ with $1\le k\le d$, and let $0\le s<k$. Then
\[
 \err_s(\chi_k)\ \ge\ 2^{-k}\sum_{j\le(k-s-1)/2}\binom kj
 \ \ge\ \max\Bigl\{2^{-k},\ \frac12-\frac{s+1}{\sqrt{2\pi k}}\Bigr\}.
\]
\end{lemma}
\begin{proof}
On the $k$-cube the first inequality is the Aspnes--Beigel--Furst--Rudich
bound \cite{ABFR}, proved in \cite[Lemma~6.2]{subcritical}. For $d>k$,
fixing the coordinates outside $[k]$ turns a degree-$s$ polynomial into a
degree-$s$ polynomial on the $k$-cube; average over the fixed
coordinates. The term $j=0$ gives $2^{-k}$. Write
$t=\lfloor(k-s-1)/2\rfloor$ and $B\sim\mathrm{Bin}(k,1/2)$. By symmetry,
$\Pp\{B\le t\}=\tfrac12(1-\Pp\{t<B<k-t\})$. The middle range contains
$k-2t-1\le s+1$ integers, each of probability at most $\sqrt{2/(\pi k)}$.
\end{proof}

\begin{corollary}[Networks without the global margin]\label{cor:signonly}
Under Assumption~\ref{ass:act}, for every depth $L\ge2$ and all gains, the
sign-only network learner satisfies, for every target $f$, every training
sequence $S$ of length $n$ and every $0\le r<d$,
\[
 \liminf_{N\to\infty}F_N(S)\ \ge\ \err_r(f)-\frac{C_0\,n}{\binom d{r+1}}.
\]
In particular:
\begin{enumerate}
\item[(a)] \emph{Full parity} ($f=\chi_d$, $d\ge25$): every training
sequence with $n\le\binom d{\lfloor\sqrt d/5\rfloor}/(6C_0)$, and hence every
one with $n\le d^{\sqrt d/10-1/2}/6000$, has $\liminf_NF_N(S)\ge1/4$.
\item[(b)] \emph{Sparse parity} ($f=\chi_k$, $25\le k\le d$): every training
sequence with $n\le\binom d{\lfloor\sqrt k/5\rfloor}/(6C_0)$ has
$\liminf_NF_N(S)\ge1/4$.
\item[(c)] \emph{Fixed $k$}: $\liminf_NF_N(S)\ge2^{-k}-C_0n/\binom dk$ for
every training sequence.
\end{enumerate}
These constants do not depend on $L$, the gains, or $\phi$.
\end{corollary}
\begin{proof}
Combine Proposition~\ref{prop:width}(iii), the inclusion
$\{q_S\le13/20\}\subseteq\{q_S<2/3\}$, and Theorem~\ref{thm:signonly}.
For (a) and (b) take $r+1=\lfloor\sqrt k/5\rfloor$ (with $k=d$ in (a)).
Lemma~\ref{lem:abfr} gives $\err_r\ge\frac12-\frac1{5\sqrt{2\pi}}>0.42$,
and the sample bound makes the subtracted term at most $1/6$. For the
explicit form in (a), $\binom dt\ge(d/t)^t$ with
$t=\lfloor\sqrt d/5\rfloor\ge\sqrt d/5-1$, so
$\log\binom dt\ge(\sqrt d/5-1)\log(5\sqrt d)\ge\frac{\sqrt d}{10}\log d
-\frac12\log d$. For (c) take $r=k-1$ and use $\err_{k-1}(\chi_k)\ge2^{-k}$.
\end{proof}

\begin{remark}[Larger readout variance]\label{rem:readout}
Corollary~\ref{cor:signonly} also covers the \emph{margin} learner when
the readout variance is $\sigma_L^2N^{-b}$ with $b<1$, all other variances
being as in \eqref{eq:weights}. This is the tuple $(0,1,\ldots,1,b)$ in the
notation of FAQ Question~9. Indeed $z_N=N^{(1-b)/2}\tilde z_N$, where
$\tilde z_N$ is the output at $b=1$. So the margin event is
$\{|\tilde z_N(x)|>\varepsilon_N\ \forall x\}$ with
$\varepsilon_N=aN^{-(1-b)/2}\to0$. For fixed $\varepsilon>0$ and large $N$,
$B_\varepsilon\subseteq B_{\varepsilon_N}\subseteq B^{\rm sgn}_S$, where
$B_\varepsilon$ adds the condition $|z|>\varepsilon$ everywhere to the
training signs. Proposition~\ref{prop:width}(i) and the portmanteau theorem
bound the liminf and limsup of $\Pp\{\tilde z_N\in B_{\varepsilon_N}\}$ by
$\Pp\{G\in B_\varepsilon\}$ and $\Pp\{G\in B^{\rm sgn}_S\}$. Letting
$\varepsilon\downarrow0$ shows that these agree in the limit, since
$\Pp\{G(x)=0\text{ for some }x\}=0$. The same holds with the extra event
$\{f(x)z(x)>0\}$. So $q_{N,S}(x)$ converges to the sign-only $q_S(x)$, as in
the $\gamma<0$ case of \cite[Proposition~2.2]{subcritical}.
\end{remark}

\begin{remark}\label{rem:signonly}
(1) The bound holds on \emph{every} training set, including deliberately
chosen ones, and uniformly in depth, gain and activation. Depth may even
grow with $d$. For comparison, \cite[Theorem~7.1]{nngpnote} proves a
$\Theta(2^d)$ bound for a different learner, the margin learner with
specified gains, at depth at least $C_a(d+\log(1/\eta)+1)$ for a fixed accuracy parameter $\eta$. (2) Theorem~\ref{thm:signonly} is
\cite[Proposition~3]{beyondparity} with the arcsine eigenvalue ratio
replaced by Lemma~\ref{lem:levels}(a). (3) For the arcsine kernel the
subcritical paper improves $\sqrt d\log d$ to $d$, for typical samples
\cite[Theorem~10.2]{subcritical}. That argument needs the sharp ratio
$\lambda_{r+1}/\tau_r\asymp1/(r\binom D{r+2})$ at $r$ linear in $d$, and
regular decay of the level masses $\binom dq\lambda_q$. For $\tanh$ the
level masses decay much faster (Section~\ref{sec:numerics}), and we have
not proved the matching estimates. The scale $\sqrt d\log d$
comes from requiring $\err_r(f)$ to stay bounded below (the
Aspnes--Beigel--Furst--Rudich lemma); with Lemma~\ref{lem:levels} alone,
larger $r$ yields at most a polynomially small failure fraction. For odd
activations the paired form improves the binomial factor, for example to
$\binom{d+1}{k+1}$ in Corollary~\ref{cor:signonly}(c) when $k$ is even,
without changing the exponent in (a) and (b).
\end{remark}

\section{Hermite coefficients of tanh and eigenvalue lower bounds}
\label{sec:tanh}
\begin{lemma}[Hermite coefficients of $\tanh$]\label{lem:tanh}
Let $\sigma>0$, $v=\E\tanh^2(\sigma G)$, and for odd $m$
\[
 I_m(\sigma)=\int_0^\infty\frac{\omega^m e^{-\sigma^2\omega^2/2}}
 {\sinh(\pi\omega/2)}\,d\omega .
\]
The dual coefficients of $\phi=\tanh$ at variance $\sigma^2$ are
$c_m(\sigma^2)=\sigma^{2m}I_m(\sigma)^2/(m!\,v)$ for odd $m$ and $0$ for
even $m$. In particular every odd coefficient is positive, and
\begin{equation}\label{eq:tanhlower}
 c_m(\sigma^2)\ \ge\ \frac{2\sqrt{2\pi}}{\sigma^2v\sqrt{m+1}}\,
 e^{-\pi\sqrt{m+1}/\sigma}\qquad(m\text{ odd}).
\end{equation}
\end{lemma}
\begin{proof}
Gaussian integration by parts gives $\E[g(G)He_m(G)]=\E g^{(m)}(G)$ for
smooth $g$ with bounded derivatives. With $g(t)=\tanh(\sigma t)$ this gives
$\beta_m=\sigma^m\E[\tanh^{(m)}(\sigma G)]/\sqrt{m!}$. The function
$\operatorname{sech}^2=\tanh'$ has Fourier transform
$\int_\R\operatorname{sech}^2(y)e^{-i\omega y}dy=\pi\omega/\sinh(\pi\omega/2)$.
This follows from the beta integral
$\int_\R e^{-i\omega y}\operatorname{sech}^2(y)dy=2\Gamma(1+\frac{i\omega}2)\Gamma(1-\frac{i\omega}2)$
(substitute $u=(1+\tanh y)/2$, so $du=\frac12\operatorname{sech}^2y\,dy$ and $e^{2y}=u/(1-u)$) and $|\Gamma(1+ix)|^2=\pi x/\sinh(\pi x)$; see \cite[Sections~3.98 and~8.33]{GR}. The transform is positive, even and exponentially decaying.
Fourier inversion and Fubini therefore give, for even $k=m-1$,
\[
 \E[\tanh^{(m)}(\sigma G)]
 =\frac1{2\pi}\int_\R\frac{\pi\omega}{\sinh(\pi\omega/2)}(i\omega)^{k}
  e^{-\sigma^2\omega^2/2}d\omega
 =(-1)^{(m-1)/2}I_m(\sigma),
\]
so $c_m=\beta_m^2/v=\sigma^{2m}I_m^2/(m!v)$. For the lower bound use
$1/\sinh y\ge2e^{-y}$. Let $J_m=\int_0^\infty\omega^me^{-\sigma^2\omega^2/2}
d\omega=\sigma^{-m-1}2^{(m-1)/2}\Gamma(\frac{m+1}2)$, and let $W$ have
density proportional to $\omega^me^{-\sigma^2\omega^2/2}$ on $(0,\infty)$.
By Jensen's inequality,
$I_m\ge2J_m\E e^{-\pi W/2}\ge2J_me^{-\pi\E W/2}$. Here
$\E W=J_{m+1}/J_m=\sqrt2\,\Gamma(\frac m2+1)/(\sigma\Gamma(\frac{m+1}2))
\le\sqrt{m+1}/\sigma$, by $\Gamma(x+\frac12)\le\sqrt x\,\Gamma(x)$. Finally
the duplication formula gives
$\sigma^{2m}J_m^2/m!=\sigma^{-2}\frac{\sqrt\pi}2\Gamma(\frac{m+1}2)/
\Gamma(\frac m2+1)\ge\sigma^{-2}\frac{\sqrt\pi}2\sqrt{2/(m+1)}$, again by
$\Gamma(x+\frac12)\le\sqrt x\,\Gamma(x)$ with $x=\frac{m+1}2$.
\end{proof}

The formula is checked numerically in \texttt{numerics/spectrum.py}: the
computed coefficients sum to one up to the truncation mass, and they
reproduce a direct quadrature of the kernel (Section~\ref{sec:numerics}).

\begin{corollary}\label{cor:composite}
For $\phi=\tanh$ and every depth and gains, the composite coefficients of
$F$ satisfy, for every odd $m$,
\[
 c_m\ \ge\ \kappa\,e^{-C_\phi\sqrt m},\qquad
 \kappa=\frac{2\sqrt{2\pi}}{\sigma_1^2v_1}\prod_{\ell=2}^{L-1}c_1(q_\ell)>0,
 \qquad C_\phi=\sqrt2\Bigl(\frac\pi{\sigma_1}+1\Bigr).
\]
\end{corollary}
\begin{proof}
Since every $\widehat\phi_{q_\ell}$ has nonnegative coefficients,
$\widehat\phi_{q_\ell}(u)\ge c_1(q_\ell)u$ coefficientwise, so
$[t^m]F\ge\prod_{\ell\ge2}c_1(q_\ell)\cdot c_m(q_1)$. Here
$c_1(q)=q(\E\operatorname{sech}^2(\sqrt qG))^2/v_\phi(q)>0$. Apply
\eqref{eq:tanhlower} with $\sigma=\sigma_1$, and use
$(m+1)^{-1/2}\ge e^{-\sqrt{m+1}}$ and $\sqrt{m+1}\le\sqrt{2m}$.
\end{proof}

\begin{lemma}[Eigenvalue lower bounds]\label{lem:eiglower}
Let $F$ be odd with $c_m\ge\kappa e^{-\psi(m)}$ for every odd $m$, where
$\kappa>0$ and $\psi$ is nondecreasing. For $d\ge101$:
\begin{enumerate}
\item[(a)] $\binom dq\lambda_q\ \ge\ \dfrac{A\kappa\,
e^{-\psi(40D)}}{2D(D+1)}$ for every $0\le q\le0.49d$;
\item[(b)] the smallest eigenvalue of the matrix $K$ satisfies
$v_{\min}(K)\ge\frac{A\kappa}8e^{-\psi(\lceil10D\log(2D)\rceil)}$.
\end{enumerate}
\end{lemma}
\begin{proof}
Poissonize the walk. If the number of steps is $\Pi_t\sim\mathrm{Poisson}(t)$, written $\Pi_t$ to avoid confusion with the width,
the coordinates flip independently, each an odd number of times with
probability $p(t)=(1-e^{-2t/D})/2$. Hence
$\sum_m\pi_t(m)p_m(j)=p(t)^j(1-p(t))^{D-j}$, where
$\pi_t(m)=e^{-t}t^m/m!$. For odd $j$, only odd $m$ contribute, and
$p_m(j)\le1$ and $\pi_t(m)\le1$. Hence, for every integer $m_1$,
\begin{equation}\label{eq:poisson}
 a_j\ \ge\ \Bigl(\min_{m\le m_1,\ m\text{ odd}}c_m\Bigr)
 \Bigl[p(t)^j(1-p(t))^{D-j}-\Pp\{\Pi_t>m_1\}\Bigr],
\end{equation}
and $\Pp\{\Pi_t>m_1\}\le(et/m_1)^{m_1}$ whenever $et<m_1$, by the Chernoff bound $\Pp\{\Pi_t\ge m\}\le e^{-t}(et/m)^m$ at $m=m_1+1$.

(a) Let $0\le q\le0.49d$ and let $j\in\{q,q+1\}$ be odd. Then
$j\le0.49d+1\le0.495D$ because $D\ge102$. Choose $t$ with
$p(t)=j/D$; then $t=\frac D2\log\frac D{D-2j}\le\frac D2\log100<2.31D$.
Take $m_1=40D$. Then $\Pp\{\Pi_t>m_1\}\le0.158^{40D}<e^{-73D}$. The
$\mathrm{Bin}(D,j/D)$ probability of its mode $j$ is at least $1/(D+1)$,
so \eqref{eq:poisson} gives
$\binom Dja_j\ge\kappa e^{-\psi(40D)}[\frac1{D+1}-2^De^{-73D}]\ge
\kappa e^{-\psi(40D)}/(2(D+1))$. Since
$\lambda_q=Aa_j$ by Lemma~\ref{lem:spectrum}(e), and
$\binom dq/\binom Dj\ge\min\{\frac{D-q}D,\frac{q+1}D\}\ge\frac1D$, part (a)
follows.

(b) Let $j\le D$ be odd, $t=\frac D2\log(2D)$ and
$m_1=\lceil10D\log(2D)\rceil\ge20t$. Then
$p(t)^j(1-p(t))^{D-j}\ge2^{-D}(1-\frac1{2D})^D\ge2^{-D-1}$, and
$\Pp\{\Pi_t>m_1\}\le(e/20)^{m_1}<2^{-D}/4$. Thus $a_j\ge\kappa
e^{-\psi(m_1)}2^{-D}/4$. Every $\lambda_q$ equals $Aa_j$ for an odd $j\le D$,
and the matrix eigenvalues are $M\lambda_q=2^{D-1}\lambda_q$.
\end{proof}

\section{The margin learner}\label{sec:margin}
Let $G\sim N(0,K)$ with $K$ as in \eqref{eq:kernel}, $F$ odd, and
\[
 A_S=\{|G(x)|>a\ \forall x;\ y_iG(X_i)>0\ \forall i\},\qquad
 \mathcal D(S,y)=-\log\Pp(A_S)
\]
for labels $y\in\{-1,1\}^n$. The posterior is the law of $G$ given $A_S$;
its relative entropy from $N(0,K)$ is exactly $\mathcal D(S,y)$. This is the
$\gamma=0$ structure of \cite[Section~7]{subcritical}, with the arcsine
kernel replaced by $K$.

\subsection{Conditioning cost}
The bump functions below are those of \cite[Lemma~7.2]{subcritical}; we
repeat the parts of that proof that involve the kernel.

\begin{lemma}[Globally margined interpolation]\label{lem:bump}
Suppose $\varepsilon_d=\min_{q\le0.49d}\binom dq\lambda_q>0$, and let
$v_{\min}=v_{\min}(K)$ (the smallest eigenvalue of the matrix $K$). Put
\[
 E_b=\frac{400}{(d+1)\varepsilon_d}+\frac{(d+1)e^{-0.0062d}}{v_{\min}} .
\]
Let $X_1,\ldots,X_n$ be independent and uniform, with $n(0.2)^{0.1d}\le0.04$.
With probability at least $1-n^2e^{-0.19d}-n\sqrt{A\,0.72^d/v_{\min}}$, the
following holds simultaneously for every labeling $y$ and every $A'>0$.
There is $h:\X\to\R$ with $|h(x)|>A'$ for all $x$, $y_ih(X_i)>A'$ for all
$i$, and
\[
 \|h\|_K^2\le8A'^2\Bigl(\frac1{\lambda_0}+8nE_b\Bigr).
\]
\end{lemma}
\begin{proof}
For $z\in\X$ let $b_z(x)=20\int_{2/5}^{9/20}(1-2p)^{\dist(x,z)}dp$. Then
$b_z(z)=1$, $0\le b_z(x)\le(1/5)^{\dist(x,z)}$, and
$\widehat{b_z}(U)=\chi_U(z)\widehat b_{|U|}$ with
$\widehat b_j=20\int_{2/5}^{9/20}(1-p)^{d-j}p^jdp$. Hence
$\binom dj\widehat b_j\le20/(d+1)$ and $\widehat b_j\le(3/5)^d$, and for
$j>0.49d$ the integrand increases in $p$, so
$\widehat b_j\le0.55^{d-j}0.45^j$ \cite[proof of Lemma~7.2]{subcritical}.
Therefore
\[
 \|b_z\|_K^2=\sum_j\binom dj\frac{\widehat b_j^2}{\lambda_j}
 \le\sum_{j\le0.49d}\frac{(20/(d+1))^2}{\varepsilon_d}
 +\frac{2^d}{v_{\min}}\sum_{j>0.49d}\binom dj0.55^{2(d-j)}0.45^{2j}
 \le E_b,
\]
using $\lambda_j\ge v_{\min}/M$. For the last sum, the exponent
$\log2+\mathsf h(t)+2(1-t)\log0.55+2t\log0.45$, with $\mathsf h$ the binary
entropy in nats, decreases for $t\ge0.49$ and is below $-0.0062$ at
$t=0.49$, as computed in \cite[proof of Lemma~7.2]{subcritical}.

Let $J(z,w)=\langle b_z,b_w\rangle_K=\sum_U\nu_{|U|}\chi_U(z)\chi_U(w)$, with
$\nu_j=\widehat b_j^2/\lambda_j\le0.36^dM/v_{\min}=0.72^d/v_{\min}$. For independent
uniform $X,Y$, Parseval gives
$\E J(X,Y)^2=\sum_j\binom dj\nu_j^2\le\nu_{\max}\|b_z\|_K^2$. Since
$E_b\ge\|b_z\|_K^2\ge b_z(z)^2/K(z,z)=1/A$, Markov's inequality gives
$\Pp\{\sum_{i\ne j}|J(X_i,X_j)|>nE_b\}\le n\sqrt{\nu_{\max}/E_b}\le
n\sqrt{A\,0.72^d/v_{\min}}$. On the complement, every subset $T$ of the indices
satisfies $\|\sum_{i\in T}b_{X_i}\|_K^2\le|T|E_b+\sum_{i\ne j}|J|\le2nE_b$.

With probability at least $1-n^2e^{-0.19d}$ all pairwise distances exceed
$0.2d$, since $\log2-\mathsf h(0.2)>0.19$. Then each $x$ has at most one
sample point within distance $0.1d$, so
$\sum_{i\in T}b_{X_i}(x)\le1+n(0.2)^{0.1d}$, and the sum is at most
$0.2+0.04$ unless $x=X_i$ with $i\in T$, where it is at least $1$. Take
$T=\{i:y_i=-1\}$ and $h=2A'(1-2\sum_{i\in T}b_{X_i})$. Then $h\le-2A'$ on
$T$ and $h\ge1.04A'$ elsewhere. Finally
$\|h\|_K^2\le8A'^2(\|1\|_K^2+4\|\sum_Tb\|_K^2)$, and
$\|1\|_K^2=1/\lambda_0$.
\end{proof}

\begin{lemma}[Conditioning cost]\label{lem:cost}
On the event of Lemma~\ref{lem:bump}, for every labeling,
\[
 \mathcal D(S,y)\le\log2+4A'^2\Bigl(\frac1{\lambda_0}+8nE_b\Bigr),
 \qquad A'=a+\sqrt{2A\log(4M)} .
\]
\end{lemma}
\begin{proof}
Let $t_0=\sqrt{2A\log(4M)}$ and $B=\{\max_x|G(x)|\le t_0\}$. A Gaussian
tail and union bound give $\Pp(B)\ge1/2$. With $h$ from
Lemma~\ref{lem:bump} at level $A'=a+t_0$, $h+B\subseteq A_S$. For a
symmetric set $B$, the Cameron--Martin formula and $\cosh\ge1$ give
$\Pp(G\in h+B)\ge\Pp(B)e^{-\|h\|_K^2/2}$, as in
\cite[proof of Lemma~7.1]{subcritical}.
\end{proof}

\subsection{Gibbs error}
\begin{lemma}[Gaussian quadratic entropy {\cite[Lemma~9.1]{subcritical}}]
\label{lem:entropy}
If $Z$ is a linear image of a centered Gaussian vector, with covariance of
trace $t$ and largest eigenvalue $\ell$, then $\E_Q\|Z\|^2\le2t+4\ell D_0$
for every law $Q$ of the underlying vector whose relative entropy from
the Gaussian law is at most $D_0$.
\end{lemma}

\begin{proposition}[Gibbs error for every target]\label{prop:gibbs}
Let $K$ be any positive-definite covariance on $\X$ that depends only on
Hamming distance and has nonincreasing degree eigenvalues. For every
target $f$, sample $S$ and $0\le s<d$, the posterior given
$A_S$ (labels $y_i=f(X_i)$) satisfies
\[
 R_S:=\E\bigl[\Pp_x\{f(x)G(x)<0\}\bigm|A_S\bigr]
 \ \ge\ \err_s(f)-\frac{2\tau_s+4\lambda_{s+1}\mathcal D(S,f)}{a^2}.
\]
\end{proposition}
\begin{proof}
Apply Lemma~\ref{lem:entropy} to $Z=M^{-1/2}G_{>s}$, whose covariance has
trace $\tau_s$ and largest eigenvalue $\lambda_{s+1}$. This gives
$\E[\|G_{>s}\|_2^2\mid A_S]\le2\tau_s+4\lambda_{s+1}\mathcal D$. Fix a
posterior draw and let $p=G_{\le s}$. Where $f(x)p(x)\le0<f(x)G(x)$ we have
$|G-p|\ge fG=|G|>a$. Hence
$\Pp_x\{fp\le0\}\le\Pp_x\{fG<0\}+\|G_{>s}\|_2^2/a^2$, and the left side is
at least $\err_s(f)$. Take the posterior expectation.
\end{proof}

\begin{theorem}[Margin learner, $\tanh$]\label{thm:margin}
Let $\phi=\tanh$, and fix $L\ge2$, gains $\sigma_1,\ldots,\sigma_L>0$ and a
margin $a>0$. There is $d_0$ such that for every $d\ge d_0$ and every
integer $0\le n\le\exp(\frac1{20}\sqrt d\log d)$ the following holds.
With probability at least $1-e^{-d/10}$ over independent uniform training
inputs, the margin learner for full parity $f=\chi_d$ satisfies
\[
 \lim_{N\to\infty}R_N(S)\ge\frac38,\qquad
 \liminf_{N\to\infty}F_N(S)\ge\frac1{16}.
\]
More generally, there is $C$, depending on $L$, the gains and $a$, such that
on the same sample event, for every target $f$ and every odd $s<d$,
\begin{equation}\label{eq:everytarget}
 \lim_{N\to\infty}R_N(S)\ \ge\ \err_s(f)-\frac{2A\,T(s)}{a^2}
 \Bigl(1+\frac{2C(d^2+n\,e^{C\sqrt d})}{\binom{d+1}{s+2}}\Bigr).
\end{equation}
\end{theorem}
\begin{proof}
\emph{Spectral inputs.} By Corollary~\ref{cor:composite},
Lemma~\ref{lem:eiglower} applies with $\psi(m)=C_\phi\sqrt m$. Thus
$\varepsilon_d\ge e^{-C\sqrt d}$ and $v_{\min}\ge e^{-C\sqrt{d\log d}}$ for large
$d$, where here and below $C$ denotes constants depending only on $L$, the
gains and $a$. In particular $E_b\le e^{C\sqrt d}$. Also
$\lambda_0=Aa_1\ge Ac_1/D$, since $p_1(1)=1/D$.

\emph{Sample event.} Let $n\le\exp(\frac1{20}\sqrt d\log d)=e^{o(d)}$. Then
$n(0.2)^{0.1d}\le0.04$, and the failure probability in
Lemma~\ref{lem:bump} is at most
$e^{-0.18d}+e^{-0.16d+C\sqrt{d\log d}}\le e^{-d/10}$ for large $d$. On this
event Lemma~\ref{lem:cost} and $A'^2\le Cd$ give, for every labeling,
$\mathcal D(S,y)\le C(d^2+n e^{C\sqrt d})$.

\emph{Every target.} Combine Proposition~\ref{prop:gibbs} with
Lemma~\ref{lem:levels}(b),(c), which give
$\lambda_{s+1}\le\tau_s/\binom D{s+2}$ and $\tau_s\le AT(s)$. Then
Proposition~\ref{prop:width}(iii) gives \eqref{eq:everytarget}.

\emph{Parity.} Let $s$ be the largest odd integer with
$s+1\le\sqrt d/5$. Lemma~\ref{lem:abfr} gives $\err_s(\chi_d)>0.42$. Since
$s+2\ge\sqrt d/5-1$,
$\log\binom D{s+2}\ge(s+2)\log\frac D{s+2}\ge\frac1{10}\sqrt d\log d-\log d$.
For $n\le\exp(\frac1{20}\sqrt d\log d)$ the fraction in
\eqref{eq:everytarget} is therefore at most
$\exp(-\frac1{20}\sqrt d\log d+C\sqrt d+C\log d)\to0$. Also $T(s)\to0$ as
$d\to\infty$, because $s\to\infty$ and $T$ does not depend on $d$
(Lemma~\ref{lem:composite}). Hence $\lim_NR_N(S)\ge0.42-o(1)\ge3/8$ for
large $d$, and \eqref{eq:mixture} gives the bound on $F_N$.
\end{proof}

\begin{remark}[Thresholds, constants, activations]\label{rem:activations}
(1) Taking $s+1\le\delta\sqrt d$ in the proof gives
$\lim_NR_N\ge\frac12-\frac\delta{\sqrt{2\pi}}-o(1)$ for
$n\le\exp(\frac\delta4\sqrt d\log d)$. For $q>1/2$ take any $\delta<\sqrt{2\pi}(q-\frac12)$ and $c_q=\delta/4$; the inequality below holds at every width and passes to the limit since $R_N\to R_S$. By the second inequality in FAQ
Question~5, $F_q\ge(R-(1-q))/q$, so for each fixed $q>1/2$ the
fraction of inputs with correct-label probability at most $q$ stays
bounded below in a range $n\le\exp(c_q\sqrt d\log d)$. (2) The constant
$\frac1{20}$ is not optimized. (3) The proof uses $\tanh$ only through
oddness, Assumption~\ref{ass:act}, and the lower bound of
Corollary~\ref{cor:composite}. It works verbatim for any odd activation
satisfying Assumption~\ref{ass:act} whose composite coefficients satisfy
$c_m\ge\kappa e^{-\psi(m)}$ for odd $m$ (by the argument of
Corollary~\ref{cor:composite} it suffices that the first-layer coefficients
satisfy this and that $c_1(q_\ell)>0$ for $\ell\ge2$), with $\psi$ nondecreasing and
$\psi(m)=o(\sqrt m\log m)$. Then $E_b=e^{o(\sqrt d\log d)}$,
$v_{\min}\ge e^{-o(d)}$, and the same threshold results. The error function is not covered: since
$\E\operatorname{erf}(\sigma U)\operatorname{erf}(\sigma V)=\frac2\pi
\arcsin\bigl(\frac{2\sigma^2\rho}{1+2\sigma^2}\bigr)$, its dual
coefficients decay geometrically in $m$. (4) For $k$-sparse
parity the same proof needs
$e^{C\sqrt d}\ll\binom D{\delta\sqrt k}$, roughly
$k\ge Cd/\log^2d$; we have not pursued smaller $k$.
\end{remark}

\section{Where the exponential proof breaks}\label{sec:obstruction}
The $\gamma=0$ part of the subcritical theorem \cite[Theorem~9.2]{subcritical}
covers $n$ up to $e^{cd}$. It chooses an adaptive odd cutoff $r$ with
$n\approx A_\delta D_r\log D_r$, where $D_r=\sum_{q\le r}\binom dq$. Its
key normalized estimate is
\[
 \frac{\mathcal D(S)}n\le Cd\rho^{-\beta},\qquad \rho=r/d,
 \qquad\text{\cite[eq.~(42)]{subcritical}},
\]
a bound polynomial in $d$ (at most $Cd^{1+\beta}$). It combines this with
the head bound
$\E^*\|G_{\le r}\|_2^2\le8\tau_r+16b_r\tau_r\mathcal D/n$
\cite[eq.~(45)]{subcritical}, where $\E^*$ is the posterior expectation
and $b_r$ bounds the sampled Gram matrix of the tail kernel. The clipped
parity lemma is then applied at the scale of the posterior output
amplitude, which is of order $d^{(1-\beta)/2}$ there
\cite[Lemma~7.5]{subcritical}. That amplitude step uses the variance
$\tau_u\asymp d^{-\beta}$ of the Fourier levels above $u\approx d/50$; at the
standard scaling this variance is tiny, and Lemma~7.5 gives no useful
amplitude. We show below that the normalization (42) also fails: the
conditioning cost per sample is superpolynomial in $d$. We then explain
what estimate would let the rest of the argument go through.

Fix $\delta\in(0,1/8)$, let $A_\delta$ be the sampling constant of
\cite[Lemma~4.4]{subcritical} (squared coercivity: with probability at
least $1-\delta$, $\frac1n\sum_i(f(X_i)P(X_i))_-^2\ge\frac14\|P\|_2^2$
for all $P$ of degree at most $r$, if $r\le d/64$ and
$n\ge A_\delta D_r\log(2D_r)$), and let $r=r(n)$ be the largest odd integer
with $A_\delta D_r\log(2D_r)\le n$.

\begin{lemma}[Universal tail-Gram bound]\label{lem:gram}
Let $F$ be odd. There are $c_1>0$ and $C_\delta$, depending only on
$\delta$ and $A_\delta$, such that for large $d$, every $n$ with
$\exp((\log(d+1))^3)\le n\le e^{c_1d}$, and $r=r(n)$, the following holds.
With probability at least $1-\delta$ over iid uniform inputs, the matrix
$Q_{SS}=((K_{>r})(X_i,X_j))_{i,j\le n}$ satisfies
\[
 \|Q_{SS}\|_{\rm op}\le b_r\tau_r,\qquad
 b_r=C_\delta\bigl(1+\sqrt{d\,rH_r}+rH_r\bigr),\qquad H_r=\log(eD/r),
\]
and $r\le d/64$, $r\ge(\log d)^2/4$.
\end{lemma}
\begin{proof}
The range of $r$ follows exactly as in \cite[Lemma~6.4]{subcritical}, which
uses only binomial estimates. That lemma's argument also gives
$n<A_\delta D_{r+2}\log(2D_{r+2})\le CA_\delta\binom D{r+2}rH_r$.
\cite[Lemma~5.2]{subcritical} bounds
$\Pp\{\|Q_{SS}-\tau I\|>z\tau\}$ in terms of four quantities. Its proof
uses only these properties of $Q$: positive semidefinite, constant diagonal
$\tau=\tau_r$, $\max_x\sum_yQ(x,y)^2=Ms_{2,r}$, $\|Q\|_{\rm op}=M\lambda_{r+1}$,
and the far-pair bound $\eta_r=\max_{d/4\le\dist\le3d/4}|Q(x,y)|$. Here
\[
 \frac{n\lambda_{r+1}}{\tau_r}\le\frac n{\binom D{r+2}}\le CA_\delta rH_r,
 \qquad
 \frac{\sqrt{n\cdot2\log M\cdot s_{2,r}}}{\tau_r}
 \le\sqrt{2d\log2\,\frac n{\binom D{r+2}}}\le C\sqrt{A_\delta d\,rH_r},
\]
by Lemma~\ref{lem:levels}(b),(d). For the far pairs, the proof of
\cite[Lemma~5.1]{subcritical} gives
$\eta_r\le\tau_r(C\sqrt De^{-r/8}+e^{-c_0D})$, once two of its steps are
reread. First, its bound for $r<j\le D-r$ uses only the Krawtchouk estimate
(20) of that paper together with $\sum_{j>r}\binom Dja_j=\tau_r/A$ (here
$Q=\sum_{j\ge r+2,\ j\text{ odd}}Aa_j\mathsf K_j$, because $F$ is odd and
$r$ is odd). Second, for $j>D-r$ it uses monotonicity:
$a_j\le a_{b}\le\tau_r/(A\binom Db)$ for an odd $b$ near $D/4$. Thus
$2\log M\cdot\eta_r/\tau_r\to0$ for $r\ge(\log d)^2/4$. The near-pair and
collision terms are at most $3ne^{d\mathsf h(1/4)}/2^d$ and
$n^2\cdot3e^{d\mathsf h(1/4)}/2^d$. They vanish when
$2c_1<\log2-\mathsf h(1/4)$. Choosing $z$ as a $\delta$-dependent multiple
of the sum of the four terms gives the claim.
\end{proof}

\begin{proposition}[The conditioning cost per sample is superpolynomial]
\label{prop:obstruction}
Let $\phi=\tanh$; fix $L$, gains, $a$ and $\delta\in(0,1/8)$. For large $d$
and every $n$ with $\exp((\log(d+1))^3)\le n\le e^{c_1d}$, with probability
at least $1-2\delta$ over iid uniform inputs labelled by full parity,
\[
 \frac{\mathcal D(S)}n\ \ge\ \frac{a^2-10\tau_r}
 {\tau_r\bigl(16b_r+4CA_\delta rH_r\bigr)} ,\qquad r=r(n).
\]
Consequently, for every $\epsilon>0$ and $k$, for all large $d$ and every
$n$ with $\exp(d^\epsilon)\le n\le e^{c_1d}$, on the same event,
$\mathcal D(S)/n\ge d^k$. In the other direction, on the event of
Lemma~\ref{lem:bump}, $\mathcal D(S)/n\le e^{C\sqrt d}$ for every $n\ge1$ with $n(0.2)^{0.1d}\le0.04$.
(Some upper limit on $n$ is necessary in the first statement: for fixed $d$,
$\mathcal D(S)$ is bounded by the cost of the full sign pattern of $f$, so
$\mathcal D(S)/n\to0$ as $n\to\infty$.)
\end{proposition}
\begin{proof}
Work on the intersection of the squared-coercivity event and the event of
Lemma~\ref{lem:gram}. Write $p=G_{\le r}$ and $\E^*$ for the posterior
expectation. Every posterior draw has $|G|>a$ everywhere, so
$\E^*\|G\|_2^2\ge a^2$. At the training inputs, $f(X_i)G(X_i)>0$ implies
$(f(X_i)p(X_i))_-\le|G_{>r}(X_i)|$. Coercivity gives
$\|p\|_2^2\le\frac4n\sum_iG_{>r}(X_i)^2$. Lemma~\ref{lem:entropy}, applied
to $n^{-1/2}(G_{>r}(X_i))_i$, whose covariance $Q_{SS}/n$ has trace
$\tau_r$ and largest eigenvalue at most $b_r\tau_r/n$, gives
$\E^*\frac1n\sum_iG_{>r}(X_i)^2\le2\tau_r+4b_r\tau_r\mathcal D/n$.
Applied to $M^{-1/2}G_{>r}$ it gives
$\E^*\|G_{>r}\|_2^2\le2\tau_r+4\lambda_{r+1}\mathcal D$, with
$\lambda_{r+1}\le\tau_r/\binom D{r+2}\le\tau_rCA_\delta rH_r/n$. Hence
\[
 a^2\le\E^*\|p\|_2^2+\E^*\|G_{>r}\|_2^2
 \le10\tau_r+\frac{\tau_r\mathcal D}n\bigl(16b_r+4CA_\delta rH_r\bigr),
\]
which is the first claim. For the second, the $\tanh$ dual activations have
all moments finite: if $h(g)=\tanh(\sqrt qg)$, the identity
$\E[h^{(k)}(G)He_j(G)]=\E[h(G)He_{j+k}(G)]$ gives
$\sum_m\frac{m!}{(m-k)!}\beta_m^2=\E[h^{(k)}(G)^2]<\infty$. Regard $(c_m)$ as a probability distribution on degrees. For a
composition $\widehat\phi\circ H$, the degree distribution is the law of
a random sum $\sum_{i\le Y}J_i$, where $Y$ has the distribution of the
coefficients of $\widehat\phi$ and the $J_i$ are independent copies of
the degree under $H$, independent of $Y$. Since
$(\sum_{i\le Y}J_i)^k\le Y^{k-1}\sum_{i\le Y}J_i^k$, we get
$\E(\sum_{i\le Y}J_i)^k\le\E Y^k\,\E J^k$. So
$T(r)\le C_kr^{-k}$ for every $k$, and $\tau_r\le AT(r)$. If
$n\ge\exp(d^\epsilon)$, then $D_{r+2}\log(2D_{r+2})>n/A_\delta$ forces
$r\ge d^{\epsilon/2}$ for large $d$, so $\tau_r\le d^{-k'}$ for every $k'$,
while $b_r+rH_r\le Cd\log d$. The upper bound follows from
Lemma~\ref{lem:cost} and $E_b\le e^{C\sqrt d}$.
\end{proof}

\paragraph{What this shows and does not show.} The proposition is about a
proof technique, not about sample complexity, and it refutes only the
normalization (42). It does not show that the subcritical strategy fails.
The clipped-parity lemma at scale $a$ only needs the posterior's
low-degree part to be $O(a)$ in $L^2$, not small, because it is multiplied
by $u/d\approx\rho$, and $\rho=r(n)/d\le\alpha$ for any fixed small $\alpha$ once $n\le e^{c_\alpha d}$ (smaller $n$ are covered by Theorem~\ref{thm:margin}). The following is an
\emph{extension suggested by the proof}, not proved here. Suppose that on
a sample event of probability at least $3/4$ there were a matching upper
bound
\begin{equation}\label{eq:matching}
 \frac{\mathcal D(S)}n\le\frac{K\,a^2}{\tau_r\,(b_r+rH_r)},
\end{equation}
with $K$ independent of $d$, $n$ and $\rho$. Work on the intersection of
that event with the coercivity event and the event of Lemma~\ref{lem:gram}
(probability at least $3/4-2\delta$), for
$\exp((\log(d+1))^3)\le n\le e^{c_1d}$, and take
$u=r+2\lceil10\log(d+1)\rceil$. The head bound would then give
$\E^*\|G_{\le r}\|_2^2\le8\tau_r+16Ka^2$. Lemma~\ref{lem:entropy} with
$\lambda_{r+1}\le\tau_r/\binom D{r+2}$ and
$n/\binom D{r+2}\le CA_\delta rH_r$ would bound the band $(r,u]$ by
$2\tau_r+CKa^2$. The same lemma with
$\lambda_{u+1}\le\tau_r/\binom D{u+2}$ and the geometric decay of
$\binom D{r+2}/\binom D{u+2}$ would make the tail above $u$ $o(1)$ (recall $\tau_r\le AT(r)\to0$, as $r\ge(\log d)^2/4$): it is at most
$2\tau_r+4Ka^2\,n/\bigl(\binom D{u+2}(b_r+rH_r)\bigr)\le2\tau_r+
CKa^2\binom D{r+2}/\binom D{u+2}$, and the binomial ratio is at most
$d^{-10}$ for $r\le\alpha d$ with $\alpha$ small, as in
\cite[Lemma~6.4]{subcritical}. The
clipped-parity lemma would then give
$|\E^*\E_xf\sign G|\le C\rho\sqrt K+o(1)$, which is an $\exp(\Omega(d))$
bound once $\rho\le\alpha$ is small enough. Proposition~\ref{prop:obstruction}
says that \eqref{eq:matching} cannot hold with a smaller right-hand side
by more than a constant factor. So the missing ingredient is an upper
bound on the conditioning cost that matches the proposition up to
constants. The only upper bound we have is the bump bound $e^{C\sqrt d}$.

A second sufficient route avoids the adaptive cutoff. It would suffice to
show that, for a small fixed $\alpha$,
\[
 \E^*\|G_{\le\alpha d}\|_2^2\le C\quad\text{and}\quad
 \E^*\|G_{>\alpha d}\|_2^2=o(1),
\]
with $C$ independent of $\alpha$. By the clipped-parity lemma at scale $a$,
$|\E_xf\sign G|\le\frac\alpha a\|G_{\le\alpha d}\|_2+\|G_{>\alpha d}\|_2/a$,
so it is enough that $\alpha\sqrt C/a\le1/4-o(1)$. The second condition
holds for $n\le e^{c_\alpha d}$ with
$c_\alpha<\min\{\mathsf h(\alpha)/2,0.09\}$, on the event of
Lemma~\ref{lem:bump}, which then has probability at least
$1-e^{-c'd}$. Indeed Lemma~\ref{lem:entropy} gives
$\E^*\|G_{>s}\|_2^2\le2\tau_s+4\lambda_{s+1}\mathcal D$, and
Lemmas~\ref{lem:levels} and~\ref{lem:cost} make the right side
$O(T(s))+e^{C\sqrt d}n/\binom D{s+2}$ at odd $s\approx\alpha d$. The first
condition is an upper bound on the posterior's low-degree amplitude, and
we do not know how to prove it. The entropy bound gives only
$\E^*\|G\|_2^2\le2A+4\lambda_0\mathcal D=O(1+\mathcal D/d)$, which is
$O(1)$ only if $\mathcal D=O(d)$; no bound here provides that.

\paragraph{Heuristic picture} (heuristic; not used above). For
$\exp(d^\epsilon)\le n\le e^{c_1d}$, the proposition and the bump
construction place the conditioning cost per sample between
$d^{\omega(1)}$ and $e^{C\sqrt d}$. The mechanism behind the lower bound is
coercivity. At the training inputs parity cannot be followed by a
low-degree polynomial, so the low-degree part is dominated by the
high-degree part there. Yet the global margin forces $\|G\|_2\ge a$. So
the posterior must place $L^2$ mass of order $a^2$ above degree $r$, where
the prior variance $\tau_r$ is tiny. If that mass sits mainly in the band
just above $r$, the cost should match the proposition, and the argument
above would give $\exp(\Omega(d))$. If instead the posterior used
large-amplitude low-degree thresholds that approximate parity, the bound
could fail. We have not ruled this out, and the numerics below are too
small to test it.

\section{Numerical illustrations}\label{sec:numerics}
All scripts and outputs are in \texttt{numerics/} next to this paper. They
are illustrations, not evidence about asymptotic exponents.

\paragraph{Spectra} (\texttt{spectrum.py}, \texttt{spectrum.txt}). We compute
the composite coefficients $c_m$ from Lemma~\ref{lem:tanh} for the unit-gain
network at $L=2,3,4$ and for the specified-gain network at $L=3$, for
$m\le600$. We then compute the cube eigenvalues through the walk
representation $a_j=\sum_mc_mp_m(j)$, which involves no cancellation, at
$d=20,40,80,160$. At $d=10$ the eigenvalues agree with a direct
Gauss--Hermite evaluation of the layer recursion followed by a Krawtchouk
transform. The checks confirm positivity, monotonicity and pairing
(Lemma~\ref{lem:spectrum}), and the inequalities of
Lemma~\ref{lem:levels}(b),(c) and \eqref{eq:tanhlower}. Over the computed
range, the level masses $\binom dq\lambda_q$ of the $\tanh$ kernels decay
much faster than those of the arcsine kernel, whose level masses decay like
$q^{-3/2}$ at $L=2$ \cite[Lemma~3.2]{subcritical}. Table~\ref{tab:spectrum} records representative values.

\begin{table}[ht]
\centering\small
\begin{tabular}{@{}lrrrrrrrr@{}}
\toprule
kernel & $d$ & $q=1$ & $q=3$ & $q=9$ & $q=17$ & $q=33$ & $q=65$ & $v_{\min}$\\
\midrule
$\tanh$ unit gain $L=2$ & 80 & -0.04 & -1.28 & -3.26 & -5.20 & -8.86 & -22.30 & -15.3 \\
$\tanh$ unit gain $L=2$ & 160 & -0.03 & -1.26 & -3.19 & -4.95 & -7.77 & -13.18 & -23.3 \\
$\tanh$ unit gain $L=3$ & 80 & -0.05 & -1.15 & -3.06 & -5.04 & -8.78 & -22.25 & -15.3 \\
$\tanh$ unit gain $L=3$ & 160 & -0.05 & -1.14 & -2.98 & -4.78 & -7.67 & -13.13 & -23.3 \\
$\tanh$ specified gain $L=3$ & 80 & -0.07 & -1.03 & -2.67 & -4.38 & -7.82 & -20.95 & -13.9 \\
$\tanh$ specified gain $L=3$ & 160 & -0.06 & -1.02 & -2.60 & -4.15 & -6.76 & -11.89 & -21.3 \\
arcsine $L=2$ & 80 & -0.20 & -0.98 & -1.74 & -2.17 & -- & -- & -1.8 \\
arcsine $L=2$ & 160 & -0.20 & -0.98 & -1.73 & -2.15 & -2.60 & -- & -2.4 \\
\bottomrule
\end{tabular}
\caption{$\log_{10}$ of the level mass $\binom dq\lambda_q/A$ of the normalized kernel $F((1+x\cdot y)/(d+1))$,
and $\log_{10}$ of the smallest eigenvalue of the matrix $F((1+x\cdot y)/(d+1))_{x,y}$. Unit gain means
$\sigma_\ell=1$ for all $\ell$; specified gain is the family of \cite[Section~7]{nngpnote} at $L=3$.
The computation truncates the slowly decaying arcsine
coefficients at $m=600$ (omitted mass $0.021$), so arcsine entries are shown only where the omitted terms
are negligible ($q\le17$ at $d=80$, $q\le33$ at $d=160$), and the arcsine $v_{\min}$ is only a lower bound. All checks listed in Section~\ref{sec:numerics} passed for every row.}
\label{tab:spectrum}
\end{table}


\paragraph{Sign-only posteriors} (\texttt{signgp\_parity.py},
\texttt{signgp\_parity.txt}). We run exact Hamiltonian Monte Carlo for the
orthant-truncated Gaussian, using the validated sampler of
\cite{beyondparity}, for full parity at $d=10,12$. The kernels are $\tanh$
unit gain at $L=2,3$, $\tanh$ specified gain at $L=3$, and the arcsine
kernel at $L=2$. Table~\ref{tab:signgp} gives the confidence-failure
fraction. In this range all four kernels fail on most unseen inputs (at least $71\%$ of them in every cell). The whole-cube fraction falls mainly because more inputs are observed. The $\tanh$ kernels fail somewhat less often than the arcsine kernel, and the three $\tanh$ kernels are nearly indistinguishable. At these dimensions the numbers say nothing about the exponent.

\begin{table}[ht]
\centering\small
\begin{tabular}{@{}lcccc@{}}
\toprule
\multicolumn{5}{@{}l}{$d=10$, $n=100,200,400,800$}\\
$\tanh$ unit gain, $L=2$ & 0.71 (0.78) & 0.62 (0.75) & 0.53 (0.79) & 0.41 (0.90)\\
$\tanh$ unit gain, $L=3$ & 0.71 (0.78) & 0.62 (0.76) & 0.54 (0.79) & 0.41 (0.90)\\
$\tanh$ specified gain, $L=3$ & 0.71 (0.79) & 0.64 (0.78) & 0.55 (0.81) & 0.42 (0.92)\\
arcsine, $L=2$ & 0.77 (0.85) & 0.75 (0.91) & 0.64 (0.94) & 0.45 (0.98)\\
\multicolumn{5}{@{}l}{$d=12$, $n=200,400,800,1600$}\\
$\tanh$ unit gain, $L=2$ & 0.71 (0.74) & 0.65 (0.71) & 0.61 (0.74) & 0.55 (0.82)\\
$\tanh$ unit gain, $L=3$ & 0.71 (0.75) & 0.65 (0.72) & 0.62 (0.75) & 0.55 (0.82)\\
$\tanh$ specified gain, $L=3$ & 0.72 (0.76) & 0.67 (0.74) & 0.63 (0.77) & 0.56 (0.83)\\
arcsine, $L=2$ & 0.82 (0.87) & 0.80 (0.88) & 0.76 (0.92) & 0.64 (0.95)\\
\bottomrule
\end{tabular}
\caption{Sign-only learner for full parity: confidence-failure fraction $F=\Pp_x\{q_S(x)\le2/3\}$ over the whole
cube (observed inputs count as successes), and in parentheses over unseen inputs only; mean over 3 seeds of
iid uniform training inputs, columns in the order of the listed $n$. Two exact-HMC chains of 1500 draws after
300 burn-in; the largest difference in $F$ between chains over all runs is 0.006.}
\label{tab:signgp}
\end{table}


\section{Status summary}
\begin{center}
\begin{tabular}{@{}p{0.45\textwidth}p{0.47\textwidth}@{}}
\toprule
Statement & Status\\
\midrule
Width limit to a nondegenerate GP, both learners & Proved (Prop.~\ref{prop:width})\\
Sign-only learner, every training set, every activation/depth/gain:
$\exp(\Omega(\sqrt d\log d))$ for full parity, $d^k$ for fixed $k$-sparse
parity & Proved (Thm.~\ref{thm:signonly}, Cor.~\ref{cor:signonly})\\
Margin learner with readout exponent $b<1$, tuple $(0,1,\ldots,1,b)$: same
bounds & Proved (Remark~\ref{rem:readout})\\
Margin learner, $\tanh$, typical samples: $\exp(\frac1{20}\sqrt d\log d)$
& Proved (Thm.~\ref{thm:margin})\\
Same for odd activations with
$c_m\ge e^{-o(\sqrt m\log m)}$ & Proved (Remark~\ref{rem:activations})\\
Per-sample conditioning cost superpolynomial for $e^{d^\epsilon}\le n\le e^{c_1d}$ (w.p.\ $\ge1-2\delta$); the subcritical normalization (42) fails & Proved (Prop.~\ref{prop:obstruction})\\
A matching upper bound on that cost would give $\exp(\Omega(d))$ & Extension suggested by the proof (Section~\ref{sec:obstruction})\\
$\exp(\Omega(d))$ for either learner at fixed depth & Open\\
Which of the two posterior behaviours in Section~\ref{sec:obstruction} occurs & Open (heuristic discussion only)\\
Margin learner for $\operatorname{erf}$, or for sparse parity with
$k\ll d/\log^2d$ & Open\\
\bottomrule
\end{tabular}
\end{center}

\begin{thebibliography}{99}
\bibitem{subcritical}
D. Vaintrob, K. H\"anni, and OpenAI Codex.
Exponential parity sample complexity below the mean-field scaling.
This collection, \texttt{main-results/subcritical-exponential}, 2026.
Lemma and equation numbers refer to its current PDF.
\bibitem{nngpnote}
D. Vaintrob, K. H\"anni, and ChatGPT.
Small circuits under Gaussian-process classification: what sample
complexity follows from the actual posterior?
This collection, \texttt{side-results/nngp-sample-complexity}, 2026.
\bibitem{beyondparity}
Claude. Which targets are hard for the subcritical learners?
This collection, \texttt{works-in-progress/beyond-parity}, 2026.
\bibitem{ABFR}
J. Aspnes, R. Beigel, M. Furst, and S. Rudich.
The expressive power of voting polynomials.
\emph{Combinatorica} 14(2):135--148, 1994.
\bibitem{GR}
I. S. Gradshteyn and I. M. Ryzhik. \emph{Table of Integrals, Series, and
Products}, 7th ed. Academic Press, 2007.
\bibitem{daniely}
A. Daniely, R. Frostig, and Y. Singer. Toward deeper understanding of
neural networks: the power of initialization and a dual view on
expressivity. \emph{NeurIPS}, 2016.
\bibitem{neal}
R. M. Neal. \emph{Bayesian Learning for Neural Networks}. Springer, 1996.
\bibitem{leeetal}
J. Lee, Y. Bahri, R. Novak, S. S. Schoenholz, J. Pennington, and
J. Sohl-Dickstein. Deep neural networks as Gaussian processes.
\emph{ICLR}, 2018.
\bibitem{matthews}
A. G. de G. Matthews, J. Hron, M. Rowland, R. E. Turner, and
Z. Ghahramani. Gaussian process behaviour in wide deep neural networks.
\emph{ICLR}, 2018.
\end{thebibliography}
\end{document}
